\section{视觉 Long CoT：把慢思考搬进视觉空间}

上一章讲语言推理，本章回到多模态。张祥雨研究完 o1 后，重新理解视觉生成可控性差的问题：很多视觉生成任务本质上也超过了单步推理复杂度。人类画图或理解空间关系时，并不是一次性从噪声生成完整图像；人会先理解场景，定位局部，检查关系，再逐步修正。多模态如果只做一次性生成，就很容易违反几何和物理约束。

\begin{figure}[H]
\centering
\includegraphics[width=0.96\textwidth]{visual-long-cot.png}
\caption{视觉空间 Long CoT：在视觉空间慢思考，把反思模式迁移到多模态推理。自制概念图，依据 01:10:57--01:24:07 对谈内容整理。}
\end{figure}

\begin{knowledgebox}{读图：视觉 CoT 不是把中文推理写长}
图中从 Image/Video 到 Visual Actions、Reason、Reflect 和 Answer。视觉 Long CoT 的重点是让模型在视觉空间做动作：局部观察、标注、回看、比较、生成辅助草稿、修正空间关系，而不只是输出更长的文字解释。
\end{knowledgebox}

\subsection{为什么生成也需要 CoT}

本节连接视觉理解和视觉生成。张祥雨的路线是先做视觉理解中的 CoT，因为理解更容易获得可验证反馈；但更远的目标是让生成也带上思维链。高可控生成不是一次性画完，而是先理解语义约束，再通过局部生成、检查、修改来逼近目标。对于复杂视频，这一点更关键，因为时间连续性、物理因果和身份一致性都需要跨帧维护。

\begin{knowledgebox}{术语消化：视觉动作空间}
\noindent\begin{tabularx}{\textwidth}{p{0.20\textwidth}p{0.34\textwidth}X}
\toprule
动作 & 含义 & 为什么有用 \\
\midrule
局部观察 & 只看图像或视频的一部分 & 降低上下文干扰，提高细节判断。 \\
标注/画框 & 在视觉空间生成中间结构 & 把隐含空间关系显式化。 \\
回看/比较 & 对不同局部或时间点做对照 & 发现不一致和物理错误。 \\
生成草稿 & 先生成可检查中间结果 & 让生成过程可反馈、可修正。 \\
反思修正 & 发现错误后重走路径 & 把视觉任务变成可搜索过程。 \\
\bottomrule
\end{tabularx}
\end{knowledgebox}

\subsection{两条腿：预训练语料与动作空间}

前面说明视觉 Long CoT 需要动作，本节把路线收束为两条腿。多模态后续路线不是单线。第一条腿是扩充预训练语料：更多高质量图文、视频、生成过程、交互过程和视觉推理数据。第二条腿是扩展动作空间：让模型不仅能看和说，还能局部观察、画图、生成、检查、调用工具、调用子模型、接收反馈。只有数据和动作空间一起扩展，视觉思考才有机会接近语言推理的跃迁。

\begin{figure}[H]
\centering
\includegraphics[width=0.94\textwidth]{two-leg-roadmap.png}
\caption{多模态后续两条腿：扩充预训练语料，同时扩展动作空间。自制概念图，依据 01:31:31--01:45:42 对谈内容整理。}
\end{figure}

\begin{knowledgebox}{读图：数据解决“见过”，动作空间解决“会做”}
左侧 Pretraining Data 让模型见到更多视觉和生成过程；右侧 Action Space 让模型在任务中有更多可执行动作。多模态 GPT-4 时刻很可能需要两者同时发生，而不是只继续堆图文对。
\end{knowledgebox}

\subsection{研究路线时间线}

本节把前面的概念压成一张时间线，帮助读者从研究史而不是热词来理解这期。张祥雨的叙述大致经历四个阶段：第一阶段是 CV scaling，把 ResNet 等模型做深做大；第二阶段是 CV 学 NLP，尝试把 Transformer、BERT/GPT、MIM 和自回归思想迁移到图像；第三阶段是多模态一体化受挫，发现生成和理解没有真正互相增强；第四阶段是 o1 之后重新理解推理，把 RL、CoT、反思和动作空间带回视觉。

\noindent\begin{tabularx}{\textwidth}{p{0.16\textwidth}p{0.31\textwidth}X}
\toprule
阶段 & 核心动作 & 得到的教训 \\
\midrule
2012--2016 & ImageNet、AlexNet、ResNet 证明视觉 model scaling & 数据、算力和模型同时扩展能打开视觉识别。 \\
2018--2022 & ViT、iGPT、MIM、对比学习等向 NLP 学习 & 架构迁移有效，但目标函数和数据属性未被根本解决。 \\
2023--2024 & 训练大多模态模型，尝试生成理解一体化 & 理解和生成各自变强，却不能自动融合。 \\
o1 之后 & 从 RL、CoT、反思和动作空间重新看多模态 & 视觉也需要慢思考和可搜索中间过程。 \\
未来两年 & long context、在线学习、自主学习和世界模型 & GPT-4 时刻可能来自环境反馈和系统协作。 \\
\bottomrule
\end{tabularx}

\begin{importantbox}{老师强调：迷茫本身是路线证据}
访谈中反复出现“做了大半年很迷茫”“后来 o1 出现后想清楚”这类表达。这里的迷茫不是噪音，而是技术路线的证据：如果一个方向通过加数据、加模型、接模块仍然没有 1+1>2，就说明缺的是新的训练目标或动作空间。
\end{importantbox}

\subsection{本章小结}

视觉 Long CoT 是这期最重要的多模态新思路：把 o1 的反思和动作空间思想迁移到视觉，让模型在视觉空间慢思考。它不是简单拉长文字链，而是给多模态模型更多可验证、可搜索、可修正的视觉动作。

